% --- Archon physics formalization source begin ---
% archon:physics
% archon:covers PhyXMiniProblems/problem_phyx_mini_0024.lean
% archon:source-report reports/phyx_mini/problem_phyx_mini_0024.source.json

\chapter{Physics problem phyx\_mini\_0024}
\label{ch:PhyXMiniProblems_problem_phyx_mini_0024}

\paragraph{Problem source.}
\#\# Physical scenario

Figure shows a top view of a square enclosure. The inner surfaces are plane mirrors. A ray of light enters a small hole in the center of one mirror.

\#\# Question

At what angle \textbackslash{}( \textbackslash{}theta \textbackslash{}) must the ray enter if it exits through the hole after being reflected once by each of the other three mirrors?

\#\# Answer choices

- A: A: \textbackslash{}( 15.0\textasciicircum{}\{\textbackslash{}circ\} \textbackslash{})

- B: B: \textbackslash{}( 60.0\textasciicircum{}\{\textbackslash{}circ\} \textbackslash{})

- C: C: \textbackslash{}( 45.0\textasciicircum{}\{\textbackslash{}circ\} \textbackslash{})

- D: D: \textbackslash{}( 30.0\textasciicircum{}\{\textbackslash{}circ\} \textbackslash{})

\#\# Figure caption (auxiliary; use the image as primary evidence)

The image shows a geometric diagram involving optics. It features a rectangular frame with what appears to be a mirror on one edge. A blue arrow, representing a light ray, enters from the bottom left, strikes the mirror, and reflects off the surface following the law of reflection.

A dashed line is drawn perpendicular to the mirror surface, representing the normal. The angle between the incoming light ray and the normal is labeled as θ (theta).

The light path includes segments indicating the path of light before and after reflection. The reflected ray follows a direction consistent with reflection over the normal line, making the angle of incidence equal to the angle of reflection.

\#\# Dataset metadata

Geometrical Optics; Spatial Relation Reasoning, Multi-Formula Reasoning

\paragraph{Recorded answer/context.}
C: C: \textbackslash{}( 45.0\textasciicircum{}\{\textbackslash{}circ\} \textbackslash{})

\paragraph{Figure/image path.}
/root/proposal\_for\_physic/science-mango/phyx\_mini\_run/phyx\_data/test\_image/24.png

\paragraph{Formalization target.}
create a compiling Lean file with sorry bodies at `PhyXMiniProblems/problem\_phyx\_mini\_0024.lean`. The Lean declarations must preserve the physical quantities, dimensions or dimensional roles, figure labels, governing-law hypotheses, and final relation expressed by this problem.
Use Mathlib/Physlib names found through LeanExplore where available. If a physics API is missing, introduce faithful local abstractions rather than scalar placeholder aliases.

\begin{theorem}[Physics formalization target]
\label{thm:physics:phyx_mini_0024:target}
\lean{PhyXMiniProblems.ProblemPhyXMini0024.entryAngle_eq_pi_div_four}
\uses{def:physics:phyx-mini-0024:phyxminiproblems-problemphyxmini0024-horizontalmirrortangent, def:physics:phyx-mini-0024:phyxminiproblems-problemphyxmini0024-verticalmirrortangent, def:physics:phyx-mini-0024:phyxminiproblems-problemphyxmini0024-specularreflectionat, def:physics:phyx-mini-0024:phyxminiproblems-problemphyxmini0024-opensegmentinside, def:physics:phyx-mini-0024:phyxminiproblems-problemphyxmini0024-onrightmirror, def:physics:phyx-mini-0024:phyxminiproblems-problemphyxmini0024-ontopmirror, def:physics:phyx-mini-0024:phyxminiproblems-problemphyxmini0024-onleftmirror, def:physics:phyx-mini-0024:phyxminiproblems-problemphyxmini0024-bottomcenteraperture, def:physics:phyx-mini-0024:phyxminiproblems-problemphyxmini0024-threemirrorrayroute, def:physics:phyx-mini-0024:phyxminiproblems-problemphyxmini0024-entryangle}
In a square mirrored enclosure, let a ray enter at the midpoint aperture of
the bottom side and then meet the right, top, and left mirrors once each
before returning to the aperture. Under specular reflection and with no
intermediate boundary contact, the entry angle from the inward normal is
$45^\circ$, answer C. Positions and lengths must be modeled in a typed
Euclidean/physical space or as explicitly named coordinate projections in one
length unit; the ray route itself is not a scalar.
\end{theorem}
\begin{proof}
Unfold the three mirror reflections into adjacent reflected copies of the
square. Equivalently, in the original square a $45^\circ$ ray from the
bottom midpoint meets successively the right midpoint, top midpoint, left
midpoint, and bottom midpoint, with equal incidence and reflection angles.
Conversely, reflecting the prescribed contact sequence into a straight line
from the original aperture to its unfolded image forces equal horizontal and
vertical displacement. Thus the direction makes $45^\circ$ with the inward
vertical normal.
\end{proof}

% --- Archon declaration topology begin ---
\section{Declaration topology}
\begin{definition}[Optical Point]
  \label{def:physics:phyx-mini-0024:phyxminiproblems-problemphyxmini0024-opticalpoint}
  \lean{PhyXMiniProblems.ProblemPhyXMini0024.OpticalPoint}
  Positions in the top-view plane of the enclosure. Space 2 is PhysLean's flat physical space with coordinates in an arbitrary, but fixed, length unit
\end{definition}
\begin{proof}
  This is the declared model definition of Optical Point.
\end{proof}
\begin{definition}[Optical Vector]
  \label{def:physics:phyx-mini-0024:phyxminiproblems-problemphyxmini0024-opticalvector}
  \lean{PhyXMiniProblems.ProblemPhyXMini0024.OpticalVector}
  Displacement and ray-direction vectors in the top-view plane
\end{definition}
\begin{proof}
  This is the declared model definition of Optical Vector.
\end{proof}
\begin{definition}[horizontal Mirror Tangent]
  \label{def:physics:phyx-mini-0024:phyxminiproblems-problemphyxmini0024-horizontalmirrortangent}
  \lean{PhyXMiniProblems.ProblemPhyXMini0024.horizontalMirrorTangent}
  \uses{def:physics:phyx-mini-0024:phyxminiproblems-problemphyxmini0024-opticalvector}
  The direction subspace tangent to either horizontal mirror of the square. Reflection of a ray direction in this subspace is specular reflection at a horizontal mirror
\end{definition}
\begin{proof}
  This is the declared model definition of horizontal Mirror Tangent.
\end{proof}
\begin{definition}[vertical Mirror Tangent]
  \label{def:physics:phyx-mini-0024:phyxminiproblems-problemphyxmini0024-verticalmirrortangent}
  \lean{PhyXMiniProblems.ProblemPhyXMini0024.verticalMirrorTangent}
  \uses{def:physics:phyx-mini-0024:phyxminiproblems-problemphyxmini0024-opticalvector}
  The direction subspace tangent to either vertical mirror of the square. Reflection of a ray direction in this subspace is specular reflection at a vertical mirror
\end{definition}
\begin{proof}
  This is the declared model definition of vertical Mirror Tangent.
\end{proof}
\begin{definition}[unit Direction]
  \label{def:physics:phyx-mini-0024:phyxminiproblems-problemphyxmini0024-unitdirection}
  \lean{PhyXMiniProblems.ProblemPhyXMini0024.unitDirection}
  \uses{def:physics:phyx-mini-0024:phyxminiproblems-problemphyxmini0024-opticalpoint, def:physics:phyx-mini-0024:phyxminiproblems-problemphyxmini0024-opticalvector}
  The unit travel direction of a nondegenerate straight ray segment
\end{definition}
\begin{proof}
  This is the declared model definition of unit Direction.
\end{proof}
\begin{definition}[Specular Reflection At]
  \label{def:physics:phyx-mini-0024:phyxminiproblems-problemphyxmini0024-specularreflectionat}
  \lean{PhyXMiniProblems.ProblemPhyXMini0024.SpecularReflectionAt}
  \uses{def:physics:phyx-mini-0024:phyxminiproblems-problemphyxmini0024-opticalpoint, def:physics:phyx-mini-0024:phyxminiproblems-problemphyxmini0024-opticalvector, def:physics:phyx-mini-0024:phyxminiproblems-problemphyxmini0024-unitdirection}
  The law of specular reflection at a plane mirror: the outgoing unit direction is the reflection of the incoming unit direction in the mirror's tangent subspace. The inequalities exclude degenerate incoming and outgoing legs
\end{definition}
\begin{proof}
  This is the declared model definition of Specular Reflection At.
\end{proof}
\begin{definition}[In Open Square]
  \label{def:physics:phyx-mini-0024:phyxminiproblems-problemphyxmini0024-inopensquare}
  \lean{PhyXMiniProblems.ProblemPhyXMini0024.InOpenSquare}
  \uses{def:physics:phyx-mini-0024:phyxminiproblems-problemphyxmini0024-opticalpoint}
  A point lies strictly inside the square enclosure. sideLength is its positive coordinate readout in the fixed length unit carried by Space 2
\end{definition}
\begin{proof}
  This is the declared model definition of In Open Square.
\end{proof}
\begin{definition}[Open Segment Inside]
  \label{def:physics:phyx-mini-0024:phyxminiproblems-problemphyxmini0024-opensegmentinside}
  \lean{PhyXMiniProblems.ProblemPhyXMini0024.OpenSegmentInside}
  \uses{def:physics:phyx-mini-0024:phyxminiproblems-problemphyxmini0024-opticalpoint, def:physics:phyx-mini-0024:phyxminiproblems-problemphyxmini0024-inopensquare}
  Except for its endpoints, a ray leg stays inside the enclosure. Thus no unlisted mirror is struck between the two consecutive labelled contacts
\end{definition}
\begin{proof}
  This is the declared model definition of Open Segment Inside.
\end{proof}
\begin{definition}[On Right Mirror]
  \label{def:physics:phyx-mini-0024:phyxminiproblems-problemphyxmini0024-onrightmirror}
  \lean{PhyXMiniProblems.ProblemPhyXMini0024.OnRightMirror}
  \uses{def:physics:phyx-mini-0024:phyxminiproblems-problemphyxmini0024-opticalpoint}
  The relative interior of the right-hand mirror in the figure
\end{definition}
\begin{proof}
  This is the declared model definition of On Right Mirror.
\end{proof}
\begin{definition}[On Top Mirror]
  \label{def:physics:phyx-mini-0024:phyxminiproblems-problemphyxmini0024-ontopmirror}
  \lean{PhyXMiniProblems.ProblemPhyXMini0024.OnTopMirror}
  \uses{def:physics:phyx-mini-0024:phyxminiproblems-problemphyxmini0024-opticalpoint}
  The relative interior of the top mirror in the figure
\end{definition}
\begin{proof}
  This is the declared model definition of On Top Mirror.
\end{proof}
\begin{definition}[On Left Mirror]
  \label{def:physics:phyx-mini-0024:phyxminiproblems-problemphyxmini0024-onleftmirror}
  \lean{PhyXMiniProblems.ProblemPhyXMini0024.OnLeftMirror}
  \uses{def:physics:phyx-mini-0024:phyxminiproblems-problemphyxmini0024-opticalpoint}
  The relative interior of the left-hand mirror in the figure
\end{definition}
\begin{proof}
  This is the declared model definition of On Left Mirror.
\end{proof}
\begin{definition}[bottom Center Aperture]
  \label{def:physics:phyx-mini-0024:phyxminiproblems-problemphyxmini0024-bottomcenteraperture}
  \lean{PhyXMiniProblems.ProblemPhyXMini0024.bottomCenterAperture}
  \uses{def:physics:phyx-mini-0024:phyxminiproblems-problemphyxmini0024-opticalpoint}
  The small aperture at the midpoint of the bottom mirror
\end{definition}
\begin{proof}
  This is the declared model definition of bottom Center Aperture.
\end{proof}
\begin{definition}[inward Normal At Aperture]
  \label{def:physics:phyx-mini-0024:phyxminiproblems-problemphyxmini0024-inwardnormalataperture}
  \lean{PhyXMiniProblems.ProblemPhyXMini0024.inwardNormalAtAperture}
  \uses{def:physics:phyx-mini-0024:phyxminiproblems-problemphyxmini0024-opticalvector}
  The inward-pointing normal to the bottom mirror at the aperture
\end{definition}
\begin{proof}
  This is the declared model definition of inward Normal At Aperture.
\end{proof}
\begin{definition}[Three Mirror Ray Route]
  \label{def:physics:phyx-mini-0024:phyxminiproblems-problemphyxmini0024-threemirrorrayroute}
  \lean{PhyXMiniProblems.ProblemPhyXMini0024.ThreeMirrorRayRoute}
  \uses{def:physics:phyx-mini-0024:phyxminiproblems-problemphyxmini0024-opticalpoint}
  The four labelled contacts of the ray route seen in the figure. The ray starts at aperture, then meets rightHit, topHit, and leftHit, before returning to aperture
\end{definition}
\begin{proof}
  This is the declared model definition of Three Mirror Ray Route.
\end{proof}
\begin{definition}[entry Angle]
  \label{def:physics:phyx-mini-0024:phyxminiproblems-problemphyxmini0024-entryangle}
  \lean{PhyXMiniProblems.ProblemPhyXMini0024.entryAngle}
  \uses{def:physics:phyx-mini-0024:phyxminiproblems-problemphyxmini0024-unitdirection, def:physics:phyx-mini-0024:phyxminiproblems-problemphyxmini0024-inwardnormalataperture, def:physics:phyx-mini-0024:phyxminiproblems-problemphyxmini0024-threemirrorrayroute}
  The entry angle θ, in radians, measured between the ray inside the enclosure and the inward normal to the mirror containing the aperture
\end{definition}
\begin{proof}
  This is the declared model definition of entry Angle.
\end{proof}
% --- Archon declaration topology end ---
% --- Archon physics formalization source end ---
